\section{整体灵敏度与鲁棒性分析}
\label{sec:overall-sensitivity}

\subsection{核数变化与继承方案的来源}
继承方案使每张图的工期随可用核数单调不增，相邻核数之间不再出现退化。表\ref{tab:best-core-choice}统计三问五核继承方案的来源核数，整图基准单列。
\begin{table}[H]\centering\small
\caption{100图五核继承方案的来源与平均加速比}\label{tab:best-core-choice}
\begin{tabular}{lrrrrrrr}\toprule
场景 & 整图基准 & 1核 & 2核 & 3核 & 4核 & 5核 & 五核平均加速比\\\midrule
A & 14 & 3 & 3 & 9 & 10 & 61 & 3.4132\\
B & 6 & 6 & 5 & 10 & 8 & 65 & 3.3862\\
C & 6 & 6 & 5 & 8 & 6 & 69 & 3.4272\\\bottomrule
\end{tabular}\end{table}

五核结果仍以五核方案为主，但A、B、C分别有14、6、6张图回到整图单核基准，另有25、29、25张图沿用1至4核方案。这些图正是固定核数方案在增核时变慢的图；继承后，四至五核的100图配对中A为61图缩短、39图不变，B为65图缩短、35图不变，C为69图缩短、31图不变。继承方案只在已评价计划之间选择，不增加评价次数。

\subsection{样本组成的重采样检验}
对100张图有放回抽样1500次，每次重新计算继承方案逐图加速比的均值，取2.5\%和97.5\%经验分位数。五核A、B、C均值分别为3.4132、3.3862、3.4272，95\%区间分别为$[3.1025,3.7151]$、$[3.0686,3.6824]$和$[3.1116,3.7245]$。固定五核计划的B/C配对也按同样方法重采样：$\log(F_B/F_C(X_B))$的均值为0.01159774，区间为$[0.00534581,0.01957431]$；C内适配比$\log(F_C(X_B)/F_C(X_C))$的均值为0.00112141，区间为$[0.00012292,0.00269912]$。两个区间均不含0，硬件效应与适配效应在图集合层面都是正的。

\subsection{统一基线与受控对照}
表\ref{tab:baseline-ablation}把相对A0的加速比与B/C同计划配对比值分列。$B(X_B^*)\to C(X_B^*)$固定计划只改变L2配置，$C(X_B^*)\to C(X_C^*)$固定C配置只改变切图与核内顺序，二者构成受控对照。
\begin{table}[H]\centering\small
\caption{五核100图的统一基线与受控对照}
\label{tab:baseline-ablation}
\setlength{\tabcolsep}{4pt}
\begin{tabularx}{\linewidth}{lYrrr}
\toprule
对照层 & 计划与配置 & 平均工期/cycle & A0加速比均值 & 配对效应比均值\\\midrule
统一基线 & 整图单Task、单核$A0$ & 3124793.73 & 1.0000 & --\\
问题一 & 继承方案$A(X_A^{\mathrm{inh}})$ & 921892.21 & 3.4132 & --\\
问题二 & 继承方案$B(X_B^{\mathrm{inh}})$ & 903275.96 & 3.3862 & --\\
问题三 & 继承方案$C(X_C^{\mathrm{inh}})$ & 895871.90 & 3.4272 & --\\
同计划无L2 & 固定核数方案$B(X_B^*)$ & 929324.47 & -- & 1.0000\\
固定计划开启L2 & $C(X_B^*)$ & 922208.33 & -- & 1.0125\\
固定L2调整计划 & $C(X_C^*)$ & 922120.10 & -- & 1.0011\\
\bottomrule
\end{tabularx}
\end{table}
平均工期是100张图完成周期的算术平均；两个比值列先逐图计算再取平均。1.0125是$F_B(X_B^*)/F_C(X_B^*)$的均值，1.0011是$F_C(X_B^*)/F_C(X_C^*)$的均值，两步合成的综合比值均值为1.0136（式\eqref{eq:cache-factors}）。

\subsection{依赖结构分组}
以第4章的原始弱连通分量数对100张图分层，比较五核继承方案的加速比，结果见表\ref{tab:topology-speedup}。分组只依据原始输入，不依据求解结果。
\begin{table}[H]\centering\small
\caption{原始计算依赖分量数与五核继承方案加速比}\label{tab:topology-speedup}
\begin{tabular}{lrrrr}\toprule
原始弱连通分量 & 图数 & A均值 & B均值 & C均值\\\midrule
1个 & 16 & 1.2963 & 1.1660 & 1.1757\\
2--100个 & 57 & 3.3607 & 3.3496 & 3.3982\\
超过100个 & 27 & 4.7783 & 4.7791 & 4.8227\\\bottomrule
\end{tabular}\end{table}

分量数与加速比明显正相关：超过100个分量的图五核均值接近4.8，而16张单分量图只有1.17至1.30。单分量图必须沿依赖切开大分量，切点处的跨核等待与搬运容易落在关键路径上，这是当前方法最主要的短板。

\subsection{候选搜索收益与评价状态}
以每个图--核数--场景的首个成功候选为起点，比较最终候选的工期，可衡量候选搜索本身的作用。问题一正式400组中117组得到更短的最终方案，平均降幅10.5562\%，五核39图改善、平均降幅16.9668\%；问题二500组中83组改善，平均降幅2.7576\%，五核18图改善、平均降幅6.7385\%；问题三500组中47组改善，平均降幅0.5006\%，五核7图改善、平均降幅0.1098\%。

候选评价的阶段结果见表\ref{tab:candidate-status}。“组合数”是1500个图--核数--场景组合，其余列以实际生成并记录的候选次数为分母，同一组合可贡献多条候选记录。
\begin{table}[H]\centering\small
\caption{主矩阵候选评价的阶段分层}\label{tab:candidate-status}
\setlength{\tabcolsep}{4pt}
\begin{tabular}{lrrrrr}\toprule
场景 & 组合数 & 成功候选 & 展开失败 & 完整评价失败 & 候选尝试数\\\midrule
A & 500 & 1000 & 214 & 0 & 1214\\
B & 500 & 890 & 0 & 110 & 1000\\
C & 500 & 894 & 0 & 77 & 971\\\midrule
合计 & 1500 & 2784 & 214 & 187 & 3185\\\bottomrule
\end{tabular}
\end{table}
候选成功率为2784/3185，1500个组合全部得到成功结果。

\subsection{评价预算灵敏度}
为考察评价次数对方案质量和计算开销的影响，另取\Case{064}、\Case{070}、\Case{093}、\Case{100}、\Case{096}和\Case{055}六张代表图，在五核配置下固定候选生成顺序，把累计完整评价预算依次设为2、5、10和20次。必评锚点先行，其余候选按轻量评分进入评价队列；成功率按该预算内已完成评价的尝试数计算，累计时间为评价耗时之和。该实验只用于观察预算趋势。
\begin{table}[H]\centering\scriptsize
\caption{六张代表图在不同评价预算下的平均结果}
\label{tab:budget-sensitivity}
\begin{tabularx}{\linewidth}{lYYYY}
\toprule
场景与指标 & 2次 & 5次 & 10次 & 20次\\\midrule
A 平均加速比 & 2.9283 & 3.2697 & 3.2697 & 3.4132\\
A 成功率 & 100.0\% & 83.3\% & 75.0\% & 71.6\%\\
A 累计评价时间/s & 0.675 & 1.211 & 1.943 & 3.969\\
B 平均加速比 & 2.9283 & 3.0916 & 3.0916 & 3.4132\\
B 成功率 & 100.0\% & 80.0\% & 65.0\% & 60.3\%\\
B 累计评价时间/s & 0.358 & 1.051 & 2.149 & 4.323\\
C 平均加速比 & 3.4166 & 3.4224 & 3.4308 & 3.4313\\
C 成功率 & 100.0\% & 100.0\% & 100.0\% & 100.0\%\\
C 累计评价时间/s & 0.363 & 0.911 & 1.812 & 3.390\\
\bottomrule
\end{tabularx}
\end{table}
注：20次预算下A、B的六图均值同为3.41315，与表\ref{tab:a-results}的100图五核均值3.4132在四位小数上恰好相同，二者样本不同，数值相同属巧合。

A场景由2次增加到20次时，六图平均加速比由2.9283升至3.4132，主要增益出现在前5次评价；C场景只从3.4166升至3.4313，累计评价时间则近似随预算线性增加。B场景在5至10次之间持平，20次时升至3.4132，但部分候选因依赖环未通过，成功率随预算下降。增加预算能改善部分图的最优工期，同时带来更多失败尝试和评价成本；对超大图应先保留整图、继承方案与首个合法切块，再按剩余时间追加评价。

\subsection{工期容差与搬运量}
在已成功评价的候选中，若允许工期不超过组内最短工期的1.005倍，正式1400组里有24组可改选搬运更少的方案，合计节省40.388 MiB；容差放宽到1.01倍时，有35组可改选，合计节省74.651 MiB。这些方案都已经过完整评价，说明小幅工期容差可以换取搬运节省；采用哪一档取决于任务对完成时间和搬运量的优先顺序。
